\begin{recipe}
[% 
    preparationtime = {\unit[10]{min}},
    portion = {\portion{1}},
    source = {},
    calory = {600kcal | 29g Protein | 93g KH | 4g Fett},
    price = {2,40€},
    created = {17.03.2026},
    %rating = {1},
]
{Fried Rice mit Soja und Gemüse}

%\graph{% pictures
%    big=pic/schnellekueche/04_fried-rice
%}

\introduction{%
Perfekt für übrig gebliebenen Reis: schneller Fried Rice mit viel Protein durch Sojageschnetzeltes. Einfach, günstig und richtig stabil nach dem Sport.
}

\ingredients{
    270 g & Gekochter Reis (am besten vom Vortag 90g roh)\index{Getreide, Pasta \& Brot!Reis}\\
    40 g & Sojageschnetzeltes\index{Tofu \& Fleischalternativen!Sojageschnetzeltes}\\
    1 & Paprika\index{Gemüse \& Pilze!Paprika}\\
    \nicefrac{1}{2} Dose & Mais\index{Vorrat \& Konserven!Mais}\\
    2--3 EL & Sojasauce\index{Öle, Saucen \& Würzmittel!Sojasauce}\\
    1--2 EL & Teriyakisauce\index{Öle, Saucen \& Würzmittel!Teriyakisauce}
}

\preparation{%
    \step%
    Sojageschnetzeltes nach Packungsanweisung in heißem Wasser einweichen und anschließend gut ausdrücken.
    
    \step%
    Paprika klein würfeln.\\
    
    \step%
    Etwas Öl in einer Pfanne erhitzen und das Sojageschnetzelte scharf anbraten.
    
    \step%
    Paprika und Mais zugeben und kurz mitbraten.\\
    
    \step%
    Reis hinzufügen und alles gut vermengen. Bei hoher Hitze anbraten, bis der Reis leicht knusprig wird.
    
    \step%
    Mit Sojasauce und Teriyakisauce abschmecken und kurz weiter braten.
}

\hint{
    Funktioniert am besten mit kaltem, trockenem Reis vom Vortag – frischer Reis wird schnell matschig.
}

\end{recipe}